% Primitive: gradient-descent — formerly GUIDE-FIRST (plan 012 D-E),
% now series-backed: plan 014 built `calculus/gradient_descent_manim.py`
% FROM this chapter (ADR 008's pipeline, book-to-screen). The scene arc
% mirrors these pages: the update derived and walked, the factor's four
% fates, the nudge-square ledger pricing the cliff, the stopping places
% stamped, the anti-ball basin hop (the series' addition — see the
% double-well truing below), and the road's walk read off one chart.
% Slotted after derivative-toolkit: everything here leans backward on
% slope-as-function and the sign-change habit. The loss walk is plan
% 010's anchor-M trajectory, reproduced independently by
% answers/gradient_descent.py; problem answers are computed by that
% script.
\chapter{Gradient descent}

\begin{narrative}
Count the road's inventory. A model's belief in a transcript is a
number the trellis can price; the logarithm turned that belief into a
loss worth minimising; softmax made the per-frame scores honest; and
the derivative toolkit reads, at any setting of any knob, which way the
loss tilts. One piece is missing between ``here is a loss'' and ``the
model trains'': how a slope becomes an \emph{update}. The answer is one
move, almost embarrassingly small, applied over and over — and this
chapter is about what the move does, what its single dial controls,
and what its stopping place does and does not certify.

\section{Walk downhill}

Let a loss $L(w)$ depend on one knob $w$. The derivative chapter's
central habit — the slope is itself a function $L'(w)$ — is the whole
apparatus needed: at the current $w$, the sign of $L'(w)$ says which
direction is uphill. So step the other way,
\[
  w \;\leftarrow\; w - \eta\, L'(w),
\]
where $\eta$, the \emph{learning rate}, converts a slope into a step
length. That single line is gradient descent.

Run it exactly once by hand, on the plainest bowl there is:
$L(w) = w^2$, so $L'(w) = 2w$. Take $\eta = \tfrac14$ and start at
$w_0 = 4$. The update is $w \leftarrow w - \tfrac12 w = \tfrac{w}{2}$:
the walk is $4 \to 2 \to 1 \to \tfrac12 \to \cdots$, halving forever.

\begin{center}
\begin{tikzpicture}[x=0.9cm, y=0.9cm]
  \draw[Muted, ->] (-3.0, 0) -- (5.4, 0) node[below] {\small $w$};
  \draw[Muted, thick] plot[domain=-2.9:4.6, samples=80] (\x, {\x*\x/4});
  \fill[GoodInk] (0, 0) circle (2.2pt);
  \node[below, color=GoodInk] at (0, -0.25) {\small slope $0$};
  % The walk 4 -> 2 -> 1 -> 1/2: chords of the bowl, halving.
  \draw[CoolInk, thick, -stealth] (4, 4) -- (2, 1);
  \draw[CoolInk, thick, -stealth] (2, 1) -- (1, 0.25);
  \draw[CoolInk, thick, -stealth] (1, 0.25) -- (0.5, 0.0625);
  \fill[AccentInk] (4, 4) circle (2.2pt);
  \fill[AccentInk] (2, 1) circle (2.2pt);
  \fill[AccentInk] (1, 0.25) circle (2.2pt);
  \fill[AccentInk] (0.5, 0.0625) circle (2.2pt);
  \node[below right, color=AccentInk] at (4.2, 3.8) {\small $w_0 = 4$};
  \node[below right, color=AccentInk] at (2.2, 0.95) {\small $w_1 = 2$};
  \node[above, color=AccentInk] at (0.75, 0.75) {\small $w_2, w_3, \dots$};
\end{tikzpicture}
\end{center}

Watch what happened without being asked for. The direction was always
downhill — the sign of the slope saw to that. And the steps shrank on
their own: near the bottom the slope is small, so $\eta\,L'(w)$ is
small. The rule brakes automatically, not because a schedule was
imposed but because the landscape flattens where it matters. On this
bowl, that is the entire story of convergence.

\section{The learning rate is a bet}

Keep the bowl, vary only $\eta$. The update is
$w \leftarrow (1 - 2\eta)\,w$ — each step multiplies the distance to
the bottom by a fixed factor:

\begin{center}
\begin{tabular}{lll}
$\eta = \tfrac14$ & factor $\tfrac12$ & glides in\\[2pt]
$\eta = \tfrac34$ & factor $-\tfrac12$ & overshoots the bottom each
step, yet converges\\[2pt]
$\eta = 1$ & factor $-1$ & ping-pongs between $4$ and $-4$ forever\\[2pt]
$\eta = \tfrac54$ & factor $-\tfrac32$ & every ``descent'' step lands
higher: diverges
\end{tabular}
\end{center}

The dial has a cliff, not a dimmer. Too small merely wastes steps; too
large converts descent into ascent, and nothing inside the rule warns
you — the update only ever reads the slope \emph{where it stands},
while the failure is a property of how sharply the landscape curves
across the whole step. The learning rate is a bet about curvature,
and the bowl collects on bad bets immediately.

\section{Where the walk stops}

The update is zero exactly where $L'(w) = 0$: gradient descent stops
at flat ground, full stop. Whether flat ground is the \emph{bottom} is
a separate question, and the derivative chapter already issued the
right habit: look at the sign change. Slope crossing $-$ to $+$ is a
valley; $+$ to $-$ is a hilltop; no crossing is a shelf. Descent
stalls at every one of them. On the double well
$L(w) = \tfrac{w^4}{4} - \tfrac{w^2}{2}$, whose slope $w^3 - w$
vanishes at $-1$, $0$, $1$, a walk started moderately to the right of
$0$ settles into the valley at $1$ (not \emph{anywhere} to the right:
at $\eta = 0.1$ a start beyond $w_0 = \sqrt{11} \approx
\anchor{014.basin.sqrt11}$ meets a
slope so steep that the first hop clears the hilltop and both valleys
— a jump the parent series makes a beat of, because no rolling ball
could copy it) — but a walk started \emph{exactly} at $0$ never
moves: it sits on the hilltop, gradient zero, perfectly stalled,
certifying nothing. A zero gradient is where the walk ends;
the sign change around it is what tells you what was found. (The
knife-edge is fragile — nudge $w_0$ to $0.1$ and the walk falls into
a valley — which is exactly why the failure is instructive rather
than common.)

\section{Many knobs, and the road's own loss}

A real model has many knobs at once, and the generalisation costs one
sentence: the \emph{gradient} is the vector collecting the loss's
slope along every knob, and the update steps against it,
$\mathbf{w} \leftarrow \mathbf{w} - \eta\,\nabla L(\mathbf{w})$. Each
knob reads its own partial slope; all move together. That — nothing
more — is \emph{plain gradient descent}, and it is the engine under
essentially all of deep
learning~\cite{probability-ian-goodfellow-yoshua-bengio-and-aaron-courville-deep-learni}.

The road can afford to watch it run on its own worked example. Make
the alignment chapter's four-frame score table trainable — each
frame's three scores driven through softmax from free knobs — and let
the loss be the negative log of the alignment sum, $-\ln P(\mathrm{AB}
\mid X)$, which starts at $\anchor{010.K.NLL}$. Plain gradient
descent at $\eta = 1$ walks it down to $\anchor{010.M.loss10}$ after
ten steps, $\anchor{010.M.loss50}$ after fifty,
$\anchor{010.M.loss200}$ after two hundred, and
$\anchor{010.M.loss5000}$ after five thousand. Read the walk's shape,
not just its endpoint: over three quarters of the loss vanished in the
first ten steps, and the last leg spent nearly five thousand steps
buying less than a hundredth. That is the parabola's lesson at scale
— flattening loss means shrinking gradient means shrinking steps —
and it is why a training curve's long, nearly level tail is normal
physiology, not a stall.

One honesty note, matching the video road's own scope: everything
real training adds — minibatch stochasticity, momentum, adaptive
per-knob steps, schedules — is refinement layered on this
one rule, and none of it appears on this road. What runs here, and
in every training curve this book shows, is the bare update, undoped.
Plain gradient descent is enough to train the road's worked example
to convergence; the refinements exist because ``many knobs'' in
practice means billions, not twelve.

\paragraph{Where this goes.} This chapter treated one object as given:
the gradient of the alignment sum itself — a loss that is a sum over
$\anchor{001.paths15}$ paths needs its slope organised as cleverly as
its value was. The next chapter computes exactly that, and the
training beats of its parent series (``plain gradient descent'',
named on screen) run precisely the walk quoted above. And the seed
sprouted: descent now has its own video series — the gradient-descent
series in `calculus/' animates this chapter beat for beat,
learning-rate cliffs, valleys and hilltops, the walk watched live,
plus one beat the screen added and these pages inherited: the
basin-hopping start no ball could copy.
\end{narrative}

\section*{Practice problems}

\begin{problem}
On $L(w) = w^2$ with $\eta = \tfrac14$ and $w_0 = 4$, write out
$w_1, w_2, w_3$ exactly. After how many steps does $|w_k|$ first drop
below $0.01$?
\begin{hint}
Substitute $L'(w) = 2w$ into the update and simplify before
iterating.
\end{hint}
\begin{solution}
The update is $w \leftarrow w - \tfrac14 \cdot 2w = \tfrac{w}{2}$, so
$w_1 = 2$, $w_2 = 1$, $w_3 = \tfrac12$, and in general
$w_k = 4 \cdot 2^{-k}$. Requiring $4 \cdot 2^{-k} < 0.01$ means
$2^k > 400$, first true at $k = 9$ (where $w_9 = 4/512 \approx
0.0078$). Nine steps — and each one cost a single multiplication,
because on this bowl the rule \emph{is} a fixed shrink factor.
\end{solution}
\end{problem}

\begin{problem}
Still on $L(w) = w^2$ from $w_0 = 4$: for
$\eta \in \{\tfrac14, \tfrac34, 1, \tfrac54\}$, find the per-step
factor and classify each walk (converges, oscillates and converges,
oscillates forever, diverges). For which $\eta$ does the walk reach
the bottom at all?
\begin{hint}
Write the update as $w \leftarrow (1 - 2\eta) w$ and ask when the
factor's magnitude is below $1$.
\end{hint}
\begin{solution}
The factors are $\tfrac12$, $-\tfrac12$, $-1$, $-\tfrac32$. With
$|1 - 2\eta| < 1$ the distance to the bottom shrinks geometrically:
$\eta = \tfrac14$ glides in, and $\eta = \tfrac34$ crosses the bottom
every step but still converges, because overshooting is fine so long
as each landing is \emph{closer}. At $\eta = 1$ the factor is $-1$:
the walk ping-pongs $4, -4, 4, \dots$ forever. At $\eta = \tfrac54$
each step lands one-and-a-half times farther out: divergence. The
walk converges exactly for $0 < \eta < 1$ — on this bowl; the
threshold is set by the bowl's curvature, which is why one landscape's
safe rate is another's cliff.
\end{solution}
\end{problem}

\begin{problem}
On the double well $L(w) = \tfrac{w^4}{4} - \tfrac{w^2}{2}$ with
$\eta = 0.1$: where does the walk settle from $w_0 = 0.5$, from
$w_0 = -0.5$, and from $w_0 = 0$? Classify each stopping place with
the sign-change test.
\begin{hint}
$L'(w) = w^3 - w$ vanishes at three points; check the slope's sign
just left and right of each.
\end{hint}
\begin{solution}
From $0.5$ the walk settles at $1$; by symmetry, from $-0.5$ it
settles at $-1$. Around $w = 1$ the slope crosses $-$ to $+$: a
valley, and likewise at $-1$. From $w_0 = 0$ the walk never moves —
$L'(0) = 0$, so the update is exactly zero — yet around $0$ the
slope crosses $+$ to $-$: a hilltop. All three are legitimate
stopping places of the same rule. Descent stops at flat ground;
only the sign change says which kind of flat ground it was.
\end{solution}
\end{problem}

\begin{problem}[stretch]
From the walk in the chapter — loss $\anchor{010.K.NLL}$ at the
start, $\anchor{010.M.loss10}$ after $10$ steps,
$\anchor{010.M.loss200}$ after $200$, $\anchor{010.M.loss5000}$
after $5000$ — compute the average per-step shrink factor of the
loss over the first ten steps, and over the last $4800$. What do the
two factors say about where a training run spends its time?
\begin{hint}
If ten steps multiply the loss by $r^{10}$, the average factor is
the tenth root of the ratio.
\end{hint}
\begin{solution}
First leg: $(\anchor{010.M.loss10} / \anchor{010.K.NLL})^{1/10}
\approx 0.86$ — the loss loses
about $14\%$ of itself \emph{every step}. Last leg:
$(\anchor{010.M.loss5000} / \anchor{010.M.loss200})^{1/4800}
\approx 0.9993$ — under a tenth of a
percent per step. The two factors differ by two orders of magnitude
in effect, which is the front-loading of the chapter in one line:
early on the gradient is large and each step earns; near the bottom
the gradient (and so the step) has shrunk, and equal-looking spans
of the loss cost geometrically more steps. Reading a training curve
means reading its slope — the long flat tail is the rule working,
not failing. (The answer script re-runs the full five-thousand-step
descent from scratch and reproduces all four anchored losses to
four decimal places.)
\end{solution}
\end{problem}
